\documentclass[11pt,a4paper]{article}
\usepackage{../shared/hanzocover}
\usepackage[utf8]{inputenc}
\usepackage[T1]{fontenc}
% Shared preamble for the compute-market series.
% One style, one vocabulary, inputted by every paper so the series reads as one voice.
\usepackage[margin=1in]{geometry}
\usepackage{booktabs}
\usepackage{amsmath}
\usepackage{amssymb}
\usepackage{amsthm}
\usepackage{graphicx}
\usepackage{xcolor}
\usepackage{hyperref}
\hypersetup{colorlinks=true, linkcolor=blue, urlcolor=blue, citecolor=blue}
\usepackage{enumitem}
\usepackage{listings}
\usepackage{array}
\usepackage{tabularx}
\usepackage{siunitx}
\sisetup{group-separator={,}, group-minimum-digits=4}

\definecolor{kw}{rgb}{0.13,0.29,0.53}
\definecolor{cmt}{rgb}{0.40,0.40,0.40}
\definecolor{str}{rgb}{0.16,0.50,0.16}
\lstset{
  basicstyle=\ttfamily\footnotesize,
  keywordstyle=\color{kw}\bfseries,
  commentstyle=\color{cmt}\itshape,
  stringstyle=\color{str},
  showstringspaces=false,
  breaklines=true,
  columns=fullflexible,
  frame=single,
  framesep=4pt,
}

\newtheorem{definition}{Definition}
\newtheorem{proposition}{Proposition}
\newtheorem{remark}{Remark}

\newcommand{\code}[1]{\texttt{#1}}
\newcommand{\repo}[1]{\texttt{#1}}

% Series vocabulary — used identically in every paper.
\newcommand{\job}{\ensuremath{J}}
\newcommand{\bundle}{\ensuremath{R}}
\newcommand{\cost}{\ensuremath{\mathcal{C}}}

\tolerance=800
\emergencystretch=3em


\title{The Price of Certainty:\\Verification Tiers and Their Cost}

\author{
Hanzo AI Research\thanks{Corresponding author: research@hanzo.ai} \\
\textit{Hanzo Industries Inc (Techstars '17)} \\
Los Angeles, California \\
\texttt{https://hanzo.ai}
}

\date{2026}

\begin{document}
\hanzocoverpage

\begin{abstract}
Verification is usually treated as a protocol constant: a system decides how
work is checked and every participant lives with that choice. We argue it should
be a purchased good with a published price curve, selected per job by the party
bearing the loss.

We enumerate six verification tiers, give each an honest cost model, and derive
the indifference conditions that separate them. The central quantity is the
ratio of value at risk to verification cost, and it varies over many orders of
magnitude across the jobs a real market will carry --- which is exactly why a
single protocol-wide soundness bar is either wasteful or unsafe, and usually
both at once for different jobs.

We also make a refinement to the claim that pricing and verification are
orthogonal. They are orthogonal in \emph{mechanism} --- the allocator need not
know whether the provider is honest, and the verifier need not know whether the
price was fair. They are coupled in the \emph{objective}, because the cost of
soundness is a term in the price. Getting this distinction right is what allows
both to be designed independently without either being wrong.
\end{abstract}

\tableofcontents
\newpage

\section{Soundness as a demand-side parameter}

A buyer posting a job faces an expected loss
\begin{equation}
\mathbb{E}[\text{loss}] \;=\; \pi \cdot (1 - \sigma) \cdot V,
\end{equation}
where $V$ is the value at risk if the answer is wrong, $\pi$ is the prior
probability that a given provider defects, and $\sigma$ is the soundness of the
verification tier --- the probability that a defection is caught. Buying tier
$T$ costs $c(T)$. The buyer's problem is
\begin{equation}
\min_{T} \; \Big[ c(T) \;+\; \pi\,(1-\sigma(T))\,V \Big],
\label{eq:tierchoice}
\end{equation}
and the optimum plainly depends on $V$, which is private to the buyer and
differs by six or more orders of magnitude between a chat completion and a
settlement instruction. No protocol constant can be right for both.

\begin{remark}
This is why we resist framing a market's verification level as a property of the
market. The market's job is to make each tier available, priced, and honestly
described. The choice belongs to whoever loses money when the answer is wrong.
\end{remark}

\section{The tiers}

\subsection{T0 --- Recomputation by the buyer}

The buyer runs the job itself. Soundness is total; cost is the job. Rational
only when the buyer wanted the answer for reasons other than lacking the
capacity to compute it --- audit sampling, spot checks at rate $q$, giving
$\sigma = q$ at cost $q \cdot c_{\text{job}}$. Random audit at low $q$ combined
with a large stake is often the cheapest adequate mechanism and is systematically
underused.

\subsection{T1 --- Byte equality and known-answer tests}

Available only for determinism class $\mathsf{D}$. The provider's output is
compared against a canonical implementation or a published vector set.
Soundness is $1$ against output corruption; cost is $O(1)$ amortised when
vectors are published, or one canonical execution otherwise. Adjudication yields
a proof of misbehaviour rather than a vote.

This tier is free in the only sense that matters --- it costs nothing beyond
what the market already publishes --- and it is unavailable to anyone who has
not done the work of making their kernels reproducible. That work is the subject
of the companion paper.

\subsection{T2 --- Freivalds and algebraic spot checks}

For the matrix products that dominate inference, correctness of $AB = C$ is
checkable without recomputing the product. Draw $r$ uniformly from
$\{0,1\}^{n}$ and test
\begin{equation}
A(Br) \;\overset{?}{=}\; Cr .
\end{equation}
If $AB \neq C$ the test fails with probability at least $1/2$; after $k$
independent rounds the soundness error is at most $2^{-k}$.

The cost is three matrix--vector products, $6n^{2}$ operations, against $2n^{3}$
for the product itself. At $n = 4096$:
\begin{equation}
\frac{2n^{3}}{6n^{2}} \;=\; \frac{n}{3} \;=\; \num{1365},
\end{equation}
so a single round costs $0.073\%$ of the product and twenty rounds --- soundness
error below $10^{-6}$ --- costs $1.5\%$. For linear-algebra-heavy work this is
the best value on the ladder by a wide margin, and it applies to class
$\mathsf{S}$ where T1 does not.

Its limits should be stated. Freivalds checks the linear algebra, not the
nonlinearities, the sampling, or that the weights were the ones paid for. It is
a strong check on the expensive part of the computation and silent about the
rest. Combining it with a weight commitment --- the provider commits to a hash of
the weight tensors, the buyer's job specification names the expected hash --- closes
the most valuable of those gaps cheaply.

\subsection{T2b --- Provenance by channel}

A case the ladder above misses entirely, and it covers a large share of real
demand: computation the buyer cannot verify because it happened inside someone
else's closed model.

Nothing algebraic helps here. There is no cheap test that a token stream came
from the model that was paid for, and no attestation is available because the
hardware belongs to a vendor who does not attest to third parties. But the
\emph{channel} is verifiable even when the computation is not. A session
terminating at the vendor's own endpoint, under a pinned certificate, with the
response carrying the vendor's model identifier, establishes provenance: the
answer came from the vendor, whatever the vendor did to produce it.

The condition is sharp and easy to get wrong. Provenance holds only for whoever
terminates the session. If the buyer holds the credential and terminates the
connection itself, the guarantee is strong and no market mechanism is needed. If
a relay holds the credential --- which is the entire point of a market --- the
relay sits between the vendor and the buyer and can fabricate the response
wholesale. Certificate pinning inside the relay proves nothing to anyone outside
it.

That gap closes by composing with the tier below: put the credential and the
pinned endpoint inside an attested boundary, and the attestation covers the relay
code while the pinned session covers the vendor. Neither half suffices; together
they give a closed-model answer a provenance chain the buyer can check.

\begin{table}[h]
\centering
\begin{tabular}{lll}
\toprule
Model class & What is verifiable & Mechanism \\
\midrule
Open weights, deterministic & the computation & T1 byte equality \\
Open weights, general & the execution & T5 attestation, T2 spot-checks \\
Closed, buyer-terminated & the source & pinned session \\
Closed, relayed & the source, conditionally & pinned session inside T5 \\
\bottomrule
\end{tabular}
\caption{One market, two verification stories. Open weights verify by execution;
closed models verify by provenance.}
\end{table}

A market serving both therefore needs both, and should say which one a given job
received rather than presenting a single undifferentiated assurance.

\subsection{T3 --- Replication}

Dispatch to $m$ independent providers, compare. With independent defection
probability $\pi$ and no collusion, undetected corruption requires all $m$ to
defect identically:
\begin{equation}
1 - \sigma \;=\; \pi^{m} \quad\text{(identical-defection model)},
\qquad c(T3) \;=\; m \cdot c_{\text{job}}.
\end{equation}
Cost is linear in $m$; soundness is exponential. The independence assumption is
the weak point and it is not a technicality: providers running the same image on
the same cloud in the same region are not independent, and a market that
allocates purely on price will systematically produce correlated committees.
Diversity has to be an explicit allocation constraint or replication silently
degrades to $m = 1$.

\subsection{T4 --- Succinct proofs}

The provider produces a proof $\pi$ that the buyer verifies in time polylogarithmic
in the computation. Soundness is cryptographic and independent of any honesty
assumption. Prover overhead relative to native execution is currently $10^{3}$
to $10^{6}$, depending on arithmetization and how badly the computation fits the
field.

Rational when the same statement is verified many times, when verification must
happen on-chain, or when $V$ is large enough that $10^{4} \times$ compute is
still cheap insurance. For a single off-chain inference it is not close.

One property matters for anything intended to outlive the current cryptographic
era: a proof system resting on pairings or discrete logarithms is Shor-breakable,
which makes it an odd foundation for a market that markets itself on quantum
readiness. Groth16, KZG and PLONK all fall in that class. Our proof substrate is
deliberately restricted to hash-based STARK/FRI with no elliptic-curve recursion
and no pairings anywhere in the pipeline, using cSHAKE256, KMAC256 and
TupleHash256 (FIPS 202 with SP 800-185) as the hash family throughout so that
Merkle commitments, FRI folding and the Fiat--Shamir transform present one
auditable surface. We note plainly that this substrate is at a stable-surface,
audit-gated stage rather than a finished one; it should be pinned to a tagged
release and treated as forward-looking rather than deployed.

\subsection{T5 --- Attested execution}

This tier has changed enough recently that treating it as a cheap curiosity, as
we did in earlier drafts of this analysis, understates it badly.

Hardware confidential computing on current accelerators encrypts device memory
and attests the device, its firmware and its configuration. The overhead for
compute-bound large-model work is small --- arithmetic dominates, and the
boundary crossings that used to cost are a minor share. Current-generation
hardware with encrypted device links removes most of the remainder. Small,
transfer-bound work still pays; the workloads a compute market actually sells
largely do not.

At a few percent overhead, attestation stops being a complement and becomes the
default for everything T1 cannot reach. That is a significant reordering: byte
equality covers bit-reproducible work and stops, while attested execution proves
that the declared computation ran on declared hardware with declared inputs ---
which is exactly the claim no algebraic check can make, and precisely the gap
that otherwise forces replication.

We classify providers by attestable capability, with all verification performed
locally rather than against a vendor's cloud service. A market that cannot verify
without an external dependency has not removed the trusted party; it has renamed
it.

\begin{table}[h]
\centering
\begin{tabular}{llll}
\toprule
Tier & Hardware & Attests & Trust band \\
\midrule
1 & accelerator-native confidential mode & device, firmware, encrypted memory & 90--100 \\
2 & confidential virtual machine plus accelerator & host measurement plus device & 70--89 \\
3 & device secure enclave with an inference engine & enclave measurement & 50--69 \\
4 & none & nothing; stake substitutes & 10--49 \\
\bottomrule
\end{tabular}
\caption{Provider classification by attestable capability. An attestation carries
device identity and model, whether confidential mode and encrypted device links
are enabled, the partitioning configuration, and driver and firmware versions.}
\end{table}

The trust root moves to a silicon vendor's certificate chain. That is a different
assumption rather than the absence of one, and the history of side-channel and
fault attacks argues against attestation as the sole mechanism at high value.
Composed with T2 --- attested execution establishing \emph{what ran}, algebraic
spot-checking establishing \emph{that the arithmetic was right} --- the pair costs
a few percent together and covers failure modes neither covers alone. For most
class-$\mathsf{S}$ volume that combination, not replication, is the right default.

\section{Choosing a tier}

Substituting into~\eqref{eq:tierchoice}, tier $T$ dominates tier $T'$ when
\begin{equation}
c(T) - c(T') \;<\; \pi \, V \, \big[\sigma(T) - \sigma(T')\big].
\end{equation}
Writing $\rho = V / c_{\text{job}}$ for the value-to-compute ratio, the
boundaries take a simple form.

\begin{table}[h]
\centering
\begin{tabular}{llll}
\toprule
Tier & Cost / $c_{\text{job}}$ & Soundness & Rational when \\
\midrule
T1 & $\approx 0$          & $1$ (class $\mathsf{D}$) & always, where available \\
T5 & $0.01$--$0.05$        & attests execution & default for class $\mathsf{S}$, pair with T2 \\
T2 & $0.001k$             & $1 - 2^{-k}$ & $\rho \gtrsim 10^{2}$, linear algebra \\
T0 & $q$                  & $q$ & $\rho \gtrsim 1/q$, with stake \\
T3 & $m$                  & $1 - \pi^{m}$ & $\rho \gtrsim 10^{3}$, diverse committee \\
T4 & $10^{3}$--$10^{6}$   & cryptographic & $\rho \gtrsim 10^{6}$, or on-chain \\
\bottomrule
\end{tabular}
\caption{Tier selection by value-to-compute ratio $\rho$.}
\end{table}

The practical reading: most volume should sit at T1 or T2, replication should be
reserved for jobs where committee diversity is genuinely achievable and the value
warrants tripling the bill, and proofs
should be reserved for statements that are settled on-chain or reused. A market
that defaults everything to replication is burning two thirds of its compute for
soundness that most of its jobs do not need.

\section{Orthogonality, stated precisely}

The claim that pricing and verification are orthogonal is right and worth
defending, but it needs a sharper statement than it usually gets.

\begin{proposition}[Mechanism orthogonality]
The allocator requires no belief about provider honesty, and the verifier
requires no belief about whether the clearing price was fair. Neither mechanism
takes the other's state as input.
\end{proposition}

\begin{proposition}[Objective coupling]
The cost of verification $c(T)$ appears as a term in the allocation objective.
A buyer requiring T3 with a diversity constraint faces a different feasible set
and a different clearing price than one requiring T1.
\end{proposition}

Both are true simultaneously, and the distinction is what makes the architecture
work: because the coupling is only through a scalar cost term, the two
mechanisms can be designed, audited, replaced and reasoned about separately.
That is a stronger and more useful statement than unqualified independence,
which would be false --- if verification were genuinely free of price
consequences, no one would ever choose a cheaper tier.

\section{What a market should publish}

Three artifacts turn this from an analysis into infrastructure. A tier registry
giving $c(T)$ and $\sigma(T)$ with stated assumptions, so buyers can compute
their own optimum rather than trusting a recommendation. A per-provider
capability declaration listing which tiers and determinism classes it honours,
which the allocator treats as a hard constraint. And a public record of caught
defections, which is what turns $\pi$ from a guess into a measurement --- without
it, every buyer is solving~\eqref{eq:tierchoice} with an invented parameter.

\end{document}
